\begin{abstract}
Extending image-based Large Multimodal Models (LMMs) to video-based LMMs always requires temporal modeling in the pre-training. However, training the temporal modules gradually erases the knowledge of visual features learned from various image-text-based scenarios, leading to degradation in some downstream tasks.
% Adapting pre-trained video-based large language models (LLMs) to downstream fine-grained video understanding tasks always requires modeling on temporal modules. However, training the temporal modules during video pretraining gradually erases the knowledge of visual features learned from various image-text-based scenarios, leading to degradation in some downstream tasks. 
% Instead of tuning video-based LLMs to downstream tasks, 
To address this issue, in this paper, we introduce a novel, efficient transfer approach termed MTransLLAMA, which employs transfer learning from pre-trained image LMMs for fine-grained video tasks with only small-scale training sets. Our method enables \textbf{fewer trainable parameters} and achieves \textbf{faster adaptation} and \textbf{higher accuracy} than pre-training video-based LMM models. Specifically, our method adopts early fusion between textual and visual features to capture fine-grained information, reuses spatial attention weights in temporal attentions for cyclical spatial-temporal reasoning, and introduces dynamic attention routing to capture both global and local information in spatial-temporal attentions. Experiments demonstrate that across multiple datasets and tasks, without relying on video pre-training, our model achieves state-of-the-art performance, enabling lightweight and efficient transfer from image-based LMMs to fine-grained video tasks.
\end{abstract}

%"The utilization of AI techniques in image-based fashion design has been the subject of growing interest in recent years. Our focus is directed towards an innovative fashion design challenge where the objective is to transfer the aesthetic of a reference image onto a clothing image while safeguarding the latter's structural integrity. This poses a formidable task as there is a lack of reference images for the newly designed fashion outputs. Although current diffusion-based image translation or neural style transfer methods allow for adaptable style transfer, it remains a challenge to retain the original image structure in a convincing manner, particularly when the reference image differs significantly from the typical appearance of clothing. To surmount this issue, we introduce a novel unsupervised structure-aware transfer approach based on diffusion modeling, which semantically creates new garments from a given clothing image and reference appearance image. The foreground clothing is separated through the utilization of automatically generated semantic masks, conditioned by labels. These masks are then employed as a guide in the denoising process to preserve the structural information. Additionally, we employ a pre-trained vision Transformer for both appearance and structure guidance. Our experiments demonstrate that our proposed method outperforms existing state-of-the-art models, producing more lifelike images in the fashion design task."